\subsection{略語}
\paragraph{表2：第10章に出てくる主な略語}
\par\smallskip\noindent\refstepcounter{table}\label{tab:ch10-02}表 \thetable: 表2：第10章に出てくる主な略語における略語・日本語・説明\par\smallskip
表 \thetable は、表2：第10章に出てくる主な略語における略語・日本語・説明を比較・確認しやすい形で整理したものである。\par\smallskip
\begin{longtable}{>{\raggedright\arraybackslash}p{0.16\linewidth}>{\raggedright\arraybackslash}p{0.30\linewidth}>{\raggedright\arraybackslash}p{0.46\linewidth}}
\toprule
略語 & 日本語 & 説明\\
\midrule
\endfirsthead
\toprule
略語 & 日本語 & 説明\\
\midrule
\endhead
BPMN & 業務プロセスモデリング表記法 & 業務や開発作業の流れを図で表すための記法である。\\
CASE & コンピュータ支援ソフトウェア工学 & 開発活動をツールで支援する考え方や環境である。\\
CMM & 能力成熟度モデル & 組織のプロセス能力や成熟度を段階的に評価するモデルである。\\
\term{CMMI}{Capability Maturity Model Integration} & 能力成熟度モデル統合 & CMM 系の考え方を統合したプロセス改善モデルである。\\
GQM & 目標・質問・メトリクス & 測定を、目標から質問、質問から測定値へ落とし込む方法である。\\
IDEF0 & 統合定義 0 & 機能や活動の入出力、制約、支援資源を表すプロセス記法である。\\
PDCA & 計画・実行・確認・改善 & プロセスを経験的に改善するための基本循環である。\\
SLCM & ソフトウェアライフサイクルモデル & 作業の並べ方や反復の仕方を示す型である。\\
SLCP & ソフトウェアライフサイクルプロセス & ライフサイクル中で実行される個々のプロセスである。\\
UML & 統一モデリング言語 & 構造や振る舞いをモデルとして表す汎用的な記法である。\\
\bottomrule
\end{longtable}
